\section{Fits with other options and stability }\label{sec:highorder}

In a phenomenological analysis, the amount of information and the accuracy one can obtain depends, obviously, on the quality of the lattice data. Given the data we have, we would like to push the limit of the analysis by including more physics in the fit until the analysis no longer yields useful information. 

\begin{figure}[tbp]
\centering
\includegraphics[width=9cm]{images/chisqmap_Eq26_clover_Pion.pdf}
\includegraphics[width=8cm]{images/Eq26_clover_Pion_Expgzmm0zdd.pdf}
\caption{
Using Eq.~(\ref{eq:logM_2}) to renormalize the $O_{\gamma_t}(z)$ matrix element in the pion state for the clover action without HYP smearing. Upper Panel: $\chi^2$ map with respect to $\lambda$ and $\Lambda_{\rm QCD}$. $\langle  \chi^{2}/{\rm d.o.f.} \rangle_{z}$ is the average of $\chi^{2}/{\rm d.o.f.}$ from fitting for each $z$. Lower Panel: renormalized matrix element for several sets of ($\Lambda_{\rm QCD}$, $\lambda$) along the 'small-$\chi^2$ band'. 
}
\label{fig:Re_overlap_quark_2}
\end{figure}

\begin{figure}[tbp]
\centering
\includegraphics[width=9cm]{images/chisqmap_Eq27_clover_Pion.pdf}
\includegraphics[width=8cm]{images/Eq27_clover_Pion_Expgzmm0zdd.pdf}
\caption{Same as Fig. \ref{fig:Re_overlap_quark_2}, fit using Eq. (\ref{eq:logM_3}).
}
\label{fig:Re_overlap_quark_3}
\end{figure}

The first question we would like to address is whether one can get more accurate information
about the linear divergence. We have used so far the one-loop perturbative
QCD prediction, but treated $\Lambda_{\rm QCD}$ as a free parameter to partially take 
into account higher order corrections. However, one might consider
directly including at least the second-order corrections. Our strategy here 
is to fix $k$ to its perturbative value, and to vary $\Lambda_{\rm QCD}$
and the explicit second order correction characterized by
a new parameter $\lambda$
\begin{align}\label{eq:logM_2}
\ln \mathcal{M}(z,a) = \frac{k' z}{a} \tilde{\alpha_{s}}(1+\lambda \tilde{\alpha_{s}})
+ g(z) + ...
\end{align}
where the terms omitted are the same as those we had before in Eq.~(\ref{eq:logM}). $k'$ is related to $k$ by $-2\pi/b_0$ (since $-k'2\pi/b_0=k$). On the other hand, one might use the $\lambda$ parameter to include
some higher order corrections~\cite{Bauer:2011ws}, 

\begin{align}\label{eq:logM_3}
\ln \mathcal{M}(z,a) = \frac{k' z}{a} \frac{\tilde{\alpha_{s}}}{1-\lambda \tilde{\alpha_{s}}}
+ g(z) + ...
\end{align}
To be consistent, we now have to use the QCD running coupling constant $\tilde{\alpha_{s}}$ up to the second order $\beta$ function $b_{1}=102-\frac{38}{3} n_{f}$,
\bea\label{eq:alpha_s2}
\tilde{\alpha_{s}}=\frac{4 \pi}{b_0 \ln (\frac{1}{a^{2} \Lambda_{\rm QCD}^{2}})}\left(1-\frac{b_1}{b_0^{2}} \times \frac{\ln[\ln (\frac{1}{a^{2} \Lambda_{\rm QCD}^{2}})]}{\ln(\frac{1}{a^{2} \Lambda_{\rm QCD}^{2}})}\right).
\eea
Therefore, again we have two global parameters $\lambda$ and $\Lambda_{\rm QCD}$ in the fit. 


The fitting process with the above parametrization is similar to what is described in Sec.~\ref{sec:reofpro}. The results are shown in Figs.~\ref{fig:Re_overlap_quark_2} and \ref{fig:Re_overlap_quark_3}. 
As one can see, the range of $\lambda$ and $\Lambda_{\rm QCD}$ are rather large, and strongly correlated. After picking several correlated values of the pair, matching to the perturbative result, the resulting residual matrix elements are shown in the lower panel. One can hardly see any difference between the different choices. Moreover, the results are essentially the same as the fit
without the extra $\lambda$ parameter (Fig.~\ref{fig:Re_overlap_quark_dff}). Therefore, we conclude
that with the data at hand, our result is stable with respect to the higher order corrections. Only with more accurate data, one might be able to see the difference of a higher-order fit. 

\begin{figure}[tbp]
\centering
\includegraphics[width=9cm]{images/chisqmap_Eq29_clover_Pion.pdf}
\includegraphics[width=8cm]{images/Eq29_clover_Pion_Expgzmm0zdd.pdf}
\caption{
Using Eq.~(\ref{eq:logM_4}) to renormalize the $O_{\gamma_t}(z)$ matrix element in the pion state for the clover action without HYP smearing. Upper Panel: $\chi^2$ map with respect to $\Lambda_{\rm QCD1}$ and $\Lambda_{\rm QCD2}$. $\langle  \chi^{2}/{\rm d.o.f.} \rangle_{z}$ is the average of $\chi^{2}/{\rm d.o.f.}$ from fitting for each $z$. Lower Panel: renormalized matrix element for several sets of ($\Lambda_{\rm QCD1}$, $\Lambda_{\rm QCD2}$) along the 'small-$\chi^2$ band'.
}
\label{fig:Re_overlap_quark_4}
\end{figure}

\begin{figure}[tbp]
\centering
\includegraphics[width=8cm]{images/clover_Pion_Expgzmm0zdd_dm.pdf}
\caption{
Renormalized $O_{\gamma_t}(z)$ correlations in the pion state for the clover action without HYP smearing based on different fitting functions: Eq.~(\ref{eq:logM}) (purple), Eq.~(\ref{eq:logM_2}) (red), Eq.~(\ref{eq:logM_3}) (blue), Eq.~(\ref{eq:logM_4}) (green).
}
\label{fig:Re_overlap_quark_dff}
\end{figure}

Another possibility one can explore is that the higher-order corrections
represented by $\Lambda_{\rm QCD}$ might be different for different ensembles. Therefore, one can discuss different  
$\Lambda_{\rm QCD}$ values for MILC and RBC data,
\begin{align}\label{eq:logM_4}
\ln \mathcal{M}(z,a) = \frac{k z}{a \ln[a \Lambda_{\rm QCD1,2}]} + g(z) + f_{1,2} (z) a \nonumber\\
+\frac{3 C_{F}}{b_0} \ln [\frac{\ln [1 /(a \Lambda_{\rm QCD1,2})]}{\ln [\mu / \Lambda_{\rm QCD1,2}]}] + \ln [1+\frac{d}{\ln (a \Lambda_{\rm QCD1,2})}],
\end{align}
where we fix the $k$ parameter and obtain a two parameter fit, $\Lambda_{\rm QCD1}$ and $\Lambda_{\rm QCD2}$. 
The results of using the above equation is shown in Fig.~\ref{fig:Re_overlap_quark_4}. As shown, the two parameters 
are strongly correlated and proportional to each other. The difference 
of the two is limited to a very small range. Furthermore, with different
choices of the correlated pair ($\Lambda_{\rm QCD1}$, $\Lambda_{\rm QCD2}$), the final renormalized results do not change. 

In Fig.~\ref{fig:Re_overlap_quark_dff}, we show the results of renormalized matrix elements based on different fitting functions. All of these fitting functions lead to the same result. Thus, given the data, it is 
unnecessary to consider other high-order terms in the fitting function and Eq.~(\ref{eq:logM}) is good enough. This also supports the argument made in Sec.~\ref{ReUncer} that we can use the fitted $\Lambda_{\rm QCD}$ in leading order to partially represent higher order corrections.

Finally, we consider $O(a^{2})$ corrections to our fitting formula. These fits were unstable, which indicates that we introduce too many fit parameters for the given data quality. The conclusion is that higher precision data is needed to constrain $O(a^2)$ terms. One can also make fits
by including  $O(a^{2})$  instead of the linear correction. The physics for doing it is
weak in our cases.
